\documentclass[12pt, fleqn, openany]{studyGuide}
\setstretch{1.4}
\enableInlineReview
\renewcommand{\VMBookTitle}{Audit Review}
\renewcommand{\VMBookSubtitle}{Grade 7 --- Ch Graphic Review --- Changed Questions}
\renewcommand{\VMFooterURL}{https://viewmath.com/grade7}
\renewcommand{\VMFooterURLDisplay}{ViewMath.com/Grade7}
\begin{document}

\begin{center}
\begin{tcolorbox}[
	enhanced,
	colback=white,
	colframe=funTeal!60,
	boxrule=1.5pt,
	arc=10pt,
	outer arc=10pt,
	width=0.9\linewidth,
	halign=center,
	top=5mm, bottom=5mm,
]
	{\Large\bfseries\sffamily\color{funTealDark}Audit Review}\\[2mm]
	{\large\sffamily\color{funGrayDark}Grade 7 --- Ch Graphic Review --- Changed Questions}\\[2mm]
	{\normalsize\sffamily\color{funGrayDark}Verify corrections below}
\end{tcolorbox}
\end{center}

\newpage
\section*{5-2-q25 \hfill \normalsize\texttt{tests\_questions\_bank\_2/topics/ch05-02-solving-two-step-equations.tex}}
\begin{practiceQuestion}{5-2-q25}{gsa}
\begin{questionText}
The balance scale is in equilibrium. Find the value of $n$.

\begin{center}
\begin{tikzpicture}[scale=0.92]
  % Fulcrum
  \draw[thick] (0,-0.5) -- (-0.4,-1.2) -- (0.4,-1.2) -- cycle;
  % Beam
  \draw[thick] (-5.4,0) -- (5.4,0);
  % Left pan
  \draw[thick] (-5.2,0) -- (-6.25,-0.82) -- (-4.15,-0.82) -- (-5.2,0);
  \draw[thick,fill=funBlue!20] (-6.25,-0.82) -- (-6.25,-1.26) -- (-4.15,-1.26) -- (-4.15,-0.82) -- cycle;
  \node at (-5.2,-1.04) {$5n - 7$};
  % Right pan
  \draw[thick] (5.2,0) -- (6.25,-0.82) -- (4.15,-0.82) -- (5.2,0);
  \draw[thick,fill=funOrange!20] (6.25,-0.82) -- (6.25,-1.26) -- (4.15,-1.26) -- (4.15,-0.82) -- cycle;
  \node at (5.2,-1.04) {$38$};
\end{tikzpicture}
\end{center}
\end{questionText}
\correctAnswer{$n = 9$}
\explanation{The balance gives $5n - 7 = 38$. Add $7$: $5n = 45$. Divide by $5$: $n = 9$. Check: $5(9) - 7 = 38$ \checkmark.}
\end{practiceQuestion}

\newpage
\section*{5-3-q24 \hfill \normalsize\texttt{tests\_questions\_bank\_2/topics/ch05-03-solving-equations-with-distributive-property.tex}}
\begin{practiceQuestion}{5-3-q24}{gmc}
\begin{questionText}
The balance scale below is in equilibrium. What is $x$?

\begin{center}
\begin{tikzpicture}[scale=0.92]
  % Fulcrum
  \draw[thick] (0,-0.5) -- (-0.4,-1.2) -- (0.4,-1.2) -- cycle;
  % Beam
  \draw[thick] (-5.4,0) -- (5.4,0);
  % Left pan: 2 groups of (x+4)
  \draw[thick] (-5.2,0) -- (-6.5,-0.82) -- (-3.9,-0.82) -- (-5.2,0);
  \draw[thick,fill=funBlue!20] (-6.5,-0.82) -- (-6.5,-1.26) -- (-3.9,-1.26) -- (-3.9,-0.82) -- cycle;
  \node at (-5.2,-1.04) {$2(x + 4)$};
  % Right pan
  \draw[thick] (5.2,0) -- (6.25,-0.82) -- (4.15,-0.82) -- (5.2,0);
  \draw[thick,fill=funOrange!20] (6.25,-0.82) -- (6.25,-1.26) -- (4.15,-1.26) -- (4.15,-0.82) -- cycle;
  \node at (5.2,-1.04) {$26$};
\end{tikzpicture}
\end{center}
\end{questionText}
\choiceA{$x = 7$}
\choiceB{$x = 9$}
\choiceC{$x = 11$}
\choiceD{$x = 13$}
\correctAnswer{B}
\explanation{The balance gives $2(x + 4) = 26$. Divide by $2$: $x + 4 = 13$. Subtract $4$: $x = 9$.}
\end{practiceQuestion}

\end{document}